\scanpage{39}
% archive sheet number: 38

1) anneau (comm. unifère) \quad $X_{\mathrm{ann}}=\widehat R$\quad
$R=\text{schémas affines de t.f.}/\mathbb Z\simeq(\text{alg de t.f. sur }\mathbb Z)^\circ$

\uline{2)} anneau local \quad $X_{\mathrm{anloc}}=\widetilde R_{\mathrm{zar}}$ \quad zar $=$ topologie de Zariski.
($U\simeq U_f\cup U_{1-f}$, cas universel $U=\mathbb A$, $f$ diagonale)

\uline{4)} anneau str. local \quad $X_{\mathrm{anstrloc}}=\widetilde R_{\mathrm{et}}$ \quad et $=$ topologie étale

(local, et pour tout $U$ et sch rel. étale $U'$ sur $U$ avec $U'\to U$ \inserted{\unsure{de schém. fini|de schém. fin.|et schémas finis}} surj.),
``surjectif'' de schémas, \deleted{\unsure{schéma|isomorphisme}} $U'$ a loc. une section sur $U$;
suffit de cas universel $U=A^n$, $U'=\operatorname{Spec\,rel}(\mathcal O_{\mathbb A^n}[T]/(F))_{F'}$,
$F(\unsure{t|T})=T^n+a_1T^{n-1}+\ldots+a_n$, les $a_i$ les $n$ projections de $A^n$, $F'$ dérivée de $F\ldots$

3) anneau local hensélien \quad $X_{\mathrm{anloc\,hen}}=\widetilde R_{\mathrm{hen}}$ \quad hen $=$ topologie hensélienne,
$X_i\to X$ est couvrant ssi $\forall$ anneau local hensélien $A$, $X_i(A)\to X(A)$ l'est

[Si $X'\to X$ dans $R$ est couvrant \inserted{$B'\leftarrow B$} au sens hens, on veut que
$\underline{\operatorname{Hom}}_{\rm an}(B',A)\to\underline{\operatorname{Hom}}_{\rm an}(B,A)$ soit épi.

\uline{5)} anneau ultralocal \quad $X_{\mathrm{anultloc}}=\widetilde R_{\mathrm{fppf}}$ \quad ?

Si $X'\to X$ \deleted{est}\inserted{dans $R$} couvrant pour fppf,
alors $A(X')\to A(X)$ est épi (il suffit de prendre $X'\to X$ \unsure{qui. fini plat|qui soit fini plat|qu. fini plat} et \unsure{même|en} fini et plat
et même dans les cas universels $U=\mathbb A^n$, $U'=\operatorname{Spec}\deleted{\unsure{A[T]/F}}\mathcal O_{\mathbb A^n}[T]/(F)$, $F$ comme dans \unsure{4}).
\uline{Question} Soit $X'\to X$ dans $R$ \deleted{tel} que $\forall$ anneau ultralocal noetherien $A$,
$X'(A)\to X(A)$ est épi, alors $X'\to X$ est-il couvrant fppf ? Résulterait du lemme : tout anneau local noeth $A$
(\unsure{ess$^t$ de|Rés. de} t.f., sur $\mathbb Z$) admet un hom local dans un anneau ultra-local qui est limite inductive d'algèbres fid. plates de t.f. sur $A$.

6) anneau réduit \quad $X_{\rm red}\simeq\widetilde R_{\deleted{\unsure{\rm top}}\rm red}\simeq\widehat R_{\rm red}$

\[
\begin{aligned}
&\underset{\unsure{\mathcal E}}{\deleted{\unsure{S}}e}
 \xrightarrow[\ ]{\sim\,0_A}\varinjlim_n
 (\underbrace{\operatorname{Ker}(A\xrightarrow{x\mapsto x^n} A)}_{\unsure{\mathrm{I\,Ens}|\mathrm{Im}|I(u)}})\\
&\Updownarrow\\
&\forall\ \deleted{\unsure{f\in}}\ V(f^n)\longrightarrow V(f)\quad \forall f\in A(U)\quad n\in\mathbb N
\end{aligned}
\]

\margin{cas universel}

$\mathrm{topred}$ la top. sur $R$ où $X_i\to Y$ est couvrant ssi $\forall U$ réduit, $X_i(U)\to Y(U)$ l'est,
i.e. s'il existe une section sur $Y_{\rm red}$ \deleted{\unsure{pour}}
$R_{\rm red}$ catégorie des \deleted{\unsure{schémas}} \inserted{objets} de $R$ réduits
(peut être \deleted{aussi} considéré comme catégorie des \inserted{sections de $R$}
\deleted{\unsure{quotients de $R$ par $B\to B_{\rm red}$ de $R$ de but réduit correspondant à un hom}} $\unsure{\leftarrow}\operatorname{Spec}$)

\uline{NB} $X_{\rm red}$ peut être défini par \uline{lim} finies
ce qui se voit sur $X_{\rm red}\simeq\widehat R_{\rm red}$ car $R_{\rm red}$ admettant des $\underset{\leftarrow}{\lim}$ finies

\uline{7)} anneau intègre\quad $X_{\rm int}\simeq\widetilde R_{\mathrm{topint}}$

a)
\diagram{\begin{tikzcd}[ampersand replacement=\&,column sep=small,row sep=small]
 A\amalg A \arrow[r,"u"]\&D\arrow[r,phantom,"\subset"]\arrow[d]\& A\times A\arrow[d,"{(x,y)\mapsto xy}"]\\
 \substack{x\mapsto(x,0)\\y\mapsto(0,y)}\&e\arrow[r,"0"]\arrow[ur,phantom,"\mathrm{cart}"]\& A
\end{tikzcd}}

\[
\begin{aligned}
 &u\ \text{\uline{un épi}}\\
 \text{b)}\qquad&\varnothing_{\unsure{\mathcal E}}\xrightarrow{\sim}\operatorname{Ker}(0_A,1_A:e\rightrightarrows A)
\end{aligned}
\]

% margin note written sideways
\margin{\unsure{a)|c)} $V(fg)=V(f)\vee V(g)$ \ $\forall f,g\in A(U)$,\quad b) $V(1_A)=\varnothing_{\mathcal E}$}

\inserted{\uline{NB} b) assuré par la famille vide couvrant l'objet initial $\varnothing_R$ de $R$, pour top int.}

top int la top. sur $R$ où $X_i\to Y$ est couvrant ssi pour toute comp. irréduct.
$Y_\alpha$ de $Y$ avec str. ind. réduite, il y a section au dessus.
\deleted{\unsure{$R_{\rm int}=$ ss-catégorie des objets intègres de $R$}}

Les objets $F\in\widetilde R_{\mathrm{top\,int}}$ sont les $F\in\widehat R$ tels que $\forall X\in\operatorname{ob}R$,
avec $X_i$, $i\in I$ \unsure{famille|famille des} comp. irréd, \unsure{ses} \deleted{\unsure{trans}}forment sommes finies en produits.
$F(X)\to\prod F(X_i)\rightrightarrows\prod F(X_i\times_X X_j)$ exact.
